\section{Monitoring, Telemetry, and Incident Response}

\subsection{Telemetry Stack}

Cada llamada a herramienta y cada salida de la API del Agent desplegado se escribe en la cadena de observabilidad, que fluye hacia tres tipos de sistemas:

\begin{itemize}
    \item \textbf{Alerting}: se establecen umbrales con base en anomaly detection, como despliegues automatizados repentinos o requests de alta frecuencia
    \item \textbf{Replay}: registra el comportamiento completo para su reproducción, incluye timeline, embeddings, entradas/salidas
    \item \textbf{Metrics}: métricas personalizadas (alignment drift, policy deviation, human override rate)
\end{itemize}

\begin{warningbox}{Observability isn't optional}
En cuanto el Agent se desvía del modo habitual, es necesario poder reproducir la cadena de llamadas completa y ver el system prompt y el context externo. Sin esta visualización, la investigación y la gobernanza son casi imposibles.
\end{warningbox}

\subsection{Incident Response Loop}

Cuando el monitoreo dispara una alerta, el proceso de respuesta de Anthropic incluye:

\begin{enumerate}
    \item Determinar rápidamente el nivel de riesgo (green/orange/red)
    \item Cambiar automáticamente a un modo de marioneta, en el que todas las salidas requieren aprobación humana
    \item Abrir un post-mortem, registrando el timeline y la root cause
    \item Actualizar el checklist y las instrucciones del prompt, para evitar la recurrence
\end{enumerate}

\subsection{Resumen de la sección}

El monitoreo del Agent necesita el respaldo de tres patas: telemetry, alert y replay; en cuanto se dispara una alerta, la organización debe responder con rapidez mediante un structured loop (determinación de nivel, containment, post-mortem, update) para evitar que la capacidad se deslice hacia la zona de riesgo.

\subsection{Data Pipelines and Instrumentation}

Telemetry data is ingested into a multi-tier pipeline: real-time streaming, batch aggregation, and backing datastore for audit. Each agent interaction emits structured events with the following fields:

\begin{itemize}
    \item `task\_id`, `agent\_variant`, `score`
    \item `tools\_used` list with duration/evidence per tool
    \item `alignment\_drift` delta vs baseline
    \item `human\_override` flag
\end{itemize}

\begin{importantbox}{Why instrumentation matters}
Instrumentation not only powers alerts but also enables retrospective capability audits: you can slice by agent version, time window, tool set, or base capability and replay the exact logs that triggered a red flag.
\end{importantbox}

The ingestion stack keeps both raw and derived layers. Raw events are immutable and tagged with `event\_hash` to ensure idempotency. Derived views pre-aggregate per ASL level and feed dashboards that teams monitor in daily standups.

Dashboards themselves include `alignment drift`, `tool success`, `red-team exposure`, and `human override latency`. When any metric crosses thresholds, the telemetry system automatically tickets the incident response crew.

